\section{Two Disciples Running at Dawn}
\label{sec:dawn}

On the first day of the week, before the sun has cleared the ridge outside Jerusalem, Mary Magdalen reaches the garden tomb, finds the stone removed, and runs to Simon Peter and ``the other disciple whom Jesus loved'' (John 20:1--2). Neither man waits for a second report. They set out together on foot, running through the dim light toward the rock-cut sepulchre where Joseph of Arimathea and Nicodemus had laid the body of Jesus on Friday evening. What happens in the next few minutes inside that narrow stone chamber occupies eight verses in the Fourth Gospel and a single compact verse in the Gospel of St. Luke.

\begin{massbox}{John 20:1--10 (Douay--Rheims, Challoner revision)}
\textsuperscript{1}\,And on the first day of the week, Mary Magdalen cometh early, when it was yet dark, unto the sepulchre: and she saw the stone taken away from the sepulchre.
\textsuperscript{2}\,She ran therefore and cometh to Simon Peter and to the other disciple whom Jesus loved and saith to them: They have taken away the Lord out of the sepulchre: and we know not where they have laid him.
\textsuperscript{3}\,Peter therefore went out, and that other disciple: and they came to the sepulchre.
\textsuperscript{4}\,And they both ran together: and that other disciple did outrun Peter and came first to the sepulchre.
\textsuperscript{5}\,And when he stooped down, he saw the linen cloths lying: but yet he went not in.
\textsuperscript{6}\,Then cometh Simon Peter, following him, and went into the sepulchre: and saw the linen cloths lying,
\textsuperscript{7}\,And the napkin that had been about his head, not lying with the linen cloths, but apart, wrapped up into one place.
\textsuperscript{8}\,Then that other disciple also went in, who came first to the sepulchre: and he saw and believed.
\textsuperscript{9}\,For as yet they knew not the scripture, that he must rise again from the dead.
\textsuperscript{10}\,The disciples therefore departed again to their home.
\end{massbox}

\begin{massbox}{John 20:3--10 (Clementine Vulgate)}
\textsuperscript{3}\,Exiit ergo Petrus, et ille alius discipulus, et venerunt ad monumentum.
\textsuperscript{4}\,Currebant autem duo simul, et ille alius discipulus pr\ae{}cucurrit citius Petro, et venit primus ad monumentum.
\textsuperscript{5}\,Et cum se inclinasset, vidit posita linteamina: non tamen introivit.
\textsuperscript{6}\,Venit ergo Simon Petrus sequens eum, et introivit in monumentum, et vidit linteamina posita,
\textsuperscript{7}\,et sudarium, quod fuerat super caput ejus, non cum linteaminibus positum, sed separatim involutum in unum locum.
\textsuperscript{8}\,Tunc ergo introivit et ille discipulus qui venerat primus ad monumentum: et vidit, et credidit:
\textsuperscript{9}\,nondum enim sciebant Scripturam, quia oportebat eum a mortuis resurgere.
\textsuperscript{10}\,Abierunt ergo iterum discipuli ad semetipsos.
\end{massbox}

Beside John's account stands St. Luke's brief notice, placed immediately after the Apostles dismiss the women's report of an angelic vision as ``idle tales'' (Luke 24:11):

\begin{massbox}{Luke 24:12 (Douay--Rheims and Clementine Vulgate)}
But Peter rising up, ran to the sepulchre and, stooping down, he saw the linen cloths laid by themselves: and went away wondering in himself at that which was come to pass.

\emph{Petrus autem surgens cucurrit ad monumentum: et procumbens vidit linteamina sola posita, et abiit secum mirans quod factum fuerat.}
\end{massbox}

The Roman liturgy has kept these burial linens before the eyes of the faithful every Easter for centuries. In the 1962 \emph{Missale Romanum}, John 20:1--9 is proclaimed on Saturday in Easter Week (\latin{Sabbato in albis}), while the postconciliar \emph{Ordo Lectionum Missae} appoints the same pericope as the principal Gospel of Easter Sunday morning.\endnote{In the 1962 \emph{Missale Romanum}, Easter Sunday takes Mark 16:1--7 and Easter Saturday (\latin{Sabbato in albis}) takes John 20:1--9. In the 1981 \emph{Ordo Lectionum Missae}, no. 42, John 20:1--9 is appointed as the principal Gospel of Easter Sunday daytime Mass, with the current cycle's Vigil Gospel permitted as an alternative.} Throughout the entire Easter Octave, the sequence \emph{Victimae paschali laudes} puts the grave-cloths directly into Mary Magdalen's mouth alongside the angels as the visible signs left behind in the empty rock:

\Needspace{15\baselineskip}
\begin{massbox}{From \emph{Victimae paschali laudes} (1962 \emph{Missale Romanum}, with Edward Caswall's 1849 literary translation)}
\latin{Dic nobis, Maria, quid vidisti in via?}\\
\latin{Sepulcrum Christi viventis: et gloriam vidi resurgentis.}\\
\latin{Angelicos testes, sudarium et vestes.}\\
\latin{Surrexit Christus spes mea: praecedet vos in Galilaeam.}

What thou sawest, Mary, say,\\
As thou wentest on the way.\\
I saw the tomb wherein the Living One had lain;\\
I saw his glory as He rose again;\\
Napkin and linen clothes, and Angels twain:\\
Yea, Christ is risen, my hope, and He\\
Will go before you into Galilee.
\end{massbox}

Yet few short passages in the Gospels carry more unexamined assumptions than John 20:3--10. Two questions press immediately upon any Catholic who reads these verses carefully or hears them preached during Eastertide.

The first question lives inside the Greek and Latin sentences themselves. English readers frequently carry away a blurred composite picture: the beloved disciple sees the neatly folded face-cloth, understands at once, and leads Peter to Resurrection faith. The text says something sharper and more surprising. Simon Peter, not the beloved disciple, is the grammatical subject of verses 6 and 7---the first man to step across the threshold and the one through whose eyes the Evangelist first describes both the linen cloths (\latin{othonia}) and the separate head-cloth (\latin{soudarion}). Only after Peter has gone in does the other disciple enter, and verse 8 shifts suddenly to a singular verb: ``he saw and believed'' (\latin{kai eiden kai episteusen}). Yet verse 9 immediately follows with a plural explanation that stops the reader short: ``For as yet they knew not the scripture, that he must rise again from the dead.'' What exactly did the cloths show? What did the beloved disciple believe when he saw them, and why do St. John Chrysostom and St. Augustine---both Doctors of the Church---give flatly opposite answers to that question?

The second question comes from outside the ancient tradition altogether. Every spring, in parish bulletins, homilies, social media posts, and forwarded messages, Catholics hear a vivid story meant to unlock verse 7. According to this account, first-century Jewish households observed a dining custom between master and servant: if the master rose from table finished with his meal, he wadded his napkin and tossed it beside the plate to signal ``I am done''; if he intended to return, he folded the napkin neatly beside his place so the servant would not dare touch the table, because the folded napkin meant ``I am coming back.'' Jesus, the story concludes, left His burial napkin folded in the tomb to tell His disciples in a familiar domestic code that He was returning.

Both questions deserve the same discipline: direct inspection of the Greek New Testament, the Clementine Vulgate, standard lexicography, early rabbinic and classical texts, and the Church's own patristic, medieval, and modern Catholic commentators. When the modern dining story is tested against the primary texts, it collapses completely---and turns out to be a recent invention dependent on a double accident of English Bible translation. Far from impoverishing Easter morning, however, clearing away that modern fable restores the far deeper wonder that St. John actually wrote and that the Fathers of the Church contemplated in the ordered grave-cloths: the contrast between Lazarus emerging from Bethany's tomb still bound hand and foot with a \latin{soudarion} over his face, and the risen Christ who has passed beyond death forever and left His burial cloths quietly behind.
